% !TeX program = xelatex
% RSI 硕士学位论文 — GB/T 7713.1-2006 结构
% 编译: latexmk -xelatex main.tex
% fontset=fandol: TeX Live 自带字体, 不扫描系统字体, 可复现
\documentclass[12pt, a4paper, openany, fontset=fandol]{ctexbook}

% ----------------------- 页面与版式 -----------------------
\usepackage[top=2.8cm, bottom=2.8cm, left=3.0cm, right=2.6cm]{geometry}
\usepackage{setspace}
\onehalfspacing
\usepackage{fancyhdr}
\pagestyle{fancy}\fancyhf{}
\fancyhead[C]{\small RSI：保真度门控的递归自改进序列推荐框架}
\fancyfoot[C]{\small\thepage}
\usepackage{emptypage}

% ----------------------- 数学与算法 -----------------------
\usepackage{amsmath, amssymb, amsthm, bm}
\usepackage{algorithm}
\usepackage{algpseudocode}

% ----------------------- 图表 -----------------------
\usepackage{graphicx}
\usepackage{tikz}
\usetikzlibrary{arrows.meta, positioning, shapes.geometric,
                fit, backgrounds, calc}
\usepackage{pgfplots}
\pgfplotsset{compat=1.18}
\usepackage{booktabs}        % 三线表
\usepackage{multirow}
\usepackage{caption}
\captionsetup{font=small, labelfont=bf}

% ----------------------- 引用与链接 -----------------------
\usepackage[numbers, sort&compress]{natbib}
\usepackage[colorlinks=true, linkcolor=black, citecolor=blue,
            urlcolor=blue]{hyperref}

% ----------------------- 计数与定理环境 -----------------------
\newtheorem{theorem}{定理}[chapter]
\newtheorem{lemma}[theorem]{引理}
\newtheorem{definition}[theorem]{定义}
\newtheorem{assumption}[theorem]{假设}
\newtheorem{proposition}[theorem]{命题}
\newtheorem{example}[theorem]{例}

% ----------------------- 符号速记 -----------------------
\newcommand{\R}{\mathbb{R}}
\newcommand{\Dset}{\mathcal{D}}
\newcommand{\Dreal}{\Dset_{\mathrm{real}}}
\newcommand{\Dsyn}{\Dset_{\mathrm{syn}}}
\newcommand{\Dmix}{\Dset_{\mathrm{mix}}^{(r)}}
\newcommand{\model}{M_{\theta}}
\newcommand{\modelr}{\model^{(r)}}
\newcommand{\dist}{\mathrm{JS}}
\newcommand{\Exp}{\mathbb{E}}
\newcommand{\grad}{\nabla_{\theta}}

% 实验结果占位宏（由 scripts/fill_numbers.py 注入真实数字）
% 实验数字注入文件（scripts/fill_numbers.py 从 results/*.json 生成）
% 论文中所有实验数字经由本文件注入，正文不出现手写数字。
\newcommand{\PopNdcg}{0.0397}
\newcommand{\PopHr}{0.0786}
\newcommand{\BprNdcg}{0.0000}
\newcommand{\BprHr}{0.0000}
\newcommand{\BaseNdcg}{0.0823}
\newcommand{\BaseHr}{0.1635}
\newcommand{\NaiveNdcg}{0.0395}
\newcommand{\NaiveHr}{0.0849}
\newcommand{\ConfNdcg}{0.0600}
\newcommand{\ConfHr}{0.1306}
\newcommand{\FixNdcg}{0.0582}
\newcommand{\FixHr}{0.1253}
\newcommand{\OursNdcg}{0.0772}
\newcommand{\OursHr}{0.1561}
\newcommand{\OursP}{nan}
\newcommand{\NaiveRoundsOne}{0.0395}
\newcommand{\NaiveFinal}{0.0395}
\newcommand{\NaiveJsPerRound}{0.000}
\newcommand{\OursDeltaText}{+0.06\%}
\newcommand{\OursMaxStep}{0.06\%}
\newcommand{\OursJsRange}{0.0003--0.0003}
\newcommand{\GateEps}{0.08}
\newcommand{\AdmCandA}{0}
\newcommand{\AdmCandB}{0}
\newcommand{\AdmCandC}{0}
\newcommand{\AdmCandD}{0}
\newcommand{\AdmJsA}{0.0000}
\newcommand{\AdmJsB}{0.0000}
\newcommand{\AdmJsC}{0.0000}
\newcommand{\AdmJsD}{0.0000}
\newcommand{\AdmPassA}{0}
\newcommand{\AdmPassB}{0}
\newcommand{\AdmPassC}{0}
\newcommand{\AdmPassD}{0}
\newcommand{\AblFull}{TBD}
\newcommand{\AblNoDist}{TBD}
\newcommand{\AblNoDistDelta}{TBD}
\newcommand{\AblFix}{TBD}
\newcommand{\AblFixDelta}{TBD}
\newcommand{\AblTopk}{TBD}
\newcommand{\AblTopkDelta}{TBD}
\newcommand{\AblNaive}{TBD}
\newcommand{\AblNaiveDelta}{TBD}
\newcommand{\AblNoDistDeltaText}{TBD}
\newcommand{\AblFixDeltaText}{TBD}
\newcommand{\AblTopkDeltaText}{TBD}
\newcommand{\GainText}{+0.0376}
\newcommand{\ExpOursEnd}{0.0772}
\newcommand{\ExpGainText}{+0.0376}
\newcommand{\SensEpsA}{TBD}
\newcommand{\SensEpsAN}{TBD}
\newcommand{\SensLamA}{TBD}
\newcommand{\SensLamAN}{TBD}
\newcommand{\SensEpsB}{TBD}
\newcommand{\SensEpsBN}{TBD}
\newcommand{\SensLamB}{TBD}
\newcommand{\SensLamBN}{TBD}
\newcommand{\SensEpsC}{TBD}
\newcommand{\SensEpsCN}{TBD}
\newcommand{\SensLamC}{TBD}
\newcommand{\SensLamCN}{TBD}
\newcommand{\SensEpsD}{TBD}
\newcommand{\SensEpsDN}{TBD}
\newcommand{\SensLamD}{TBD}
\newcommand{\SensLamDN}{TBD}
\newcommand{\SensEpsE}{TBD}
\newcommand{\SensEpsEN}{TBD}
\newcommand{\SensLamE}{TBD}
\newcommand{\SensLamEN}{TBD}
\newcommand{\SprSparse}{TBD}
\newcommand{\SprSparseN}{TBD}
\newcommand{\SprMedium}{TBD}
\newcommand{\SprMediumN}{TBD}
\newcommand{\SprDense}{TBD}
\newcommand{\SprDenseN}{TBD}
\newcommand{\ExpNaiveStart}{0.0766}
\newcommand{\ExpNaiveEnd}{\NaiveFinal}
\newcommand{\ExpNaiveRounds}{4}
\newcommand{\ExpRounds}{4}


\begin{document}

% ======================= 前置部分 =======================
\frontmatter
% 封面（占位信息由作者按学校模板替换）
\begin{titlepage}
\begin{center}
\vspace*{1.5cm}

{\zihao{4}\heiti 硕 士 学 位 论 文}\\[1.2cm]

{\zihao{0}\heiti 保真度门控的递归自改进序列推荐框架研究}\\[0.8cm]

{\zihao{-2} Fidelity-Gated Recursive Self-Improvement\\for Sequential Recommendation}\\[2.5cm]

\begin{tabular}{rl}
\zihao{-4}\heiti 学\hspace{2em}号： & \underline{\makebox[7cm]{〈填写学号〉}}\\[0.8em]
\heiti 姓\hspace{2em}名： & \underline{\makebox[7cm]{〈填写姓名〉}}\\[0.8em]
\heiti 专\hspace{2em}业： & \underline{\makebox[7cm]{计算机科学与技术}}\\[0.8em]
\heiti 研究方向： & \underline{\makebox[7cm]{递归自改进 / 推荐系统自迭代}}\\[0.8em]
\heiti 指导教师： & \underline{\makebox[7cm]{〈填写导师〉}}\\
\end{tabular}

\vspace{2.5cm}

{\zihao{4} 二〇二六年六月}\\
\end{center}
\end{titlepage}

% 原创性声明与版权授权（按学校模板替换）
\chapter*{原创性声明}
\thispagestyle{empty}
本人郑重声明：所呈交的学位论文，是本人在导师指导下，独立进行研究
工作所取得的成果。除文中已经注明引用的内容外，本论文不包含任何
其他个人或集体已经发表或撰写过的作品或成果。对本文的研究作出重
要贡献的个人和集体，均已在文中以明确方式标明。本声明的法律责任
由本人承担。

\vspace{2em}
\hfill 论文作者签名：\underline{\makebox[4cm]{}} \quad 日期：\underline{\makebox[2.5cm]{}}

\vspace{4em}

\chapter*{学位论文版权使用授权书}
\thispagestyle{empty}
本学位论文作者及指导教师完全了解学校有关保留、使用学位论文的
规定，同意学校保留并向国家有关部门或机构送交论文的复印件和电子
版，允许论文被查阅和借阅。本人授权学校可以将本学位论文的全部或
部分内容编入有关数据库进行检索，可以采用影印、缩印或扫描等复制
手段保存和汇编本学位论文。

\hfill 作者签名：\underline{\makebox[3cm]{}} \quad 导师签名：\underline{\makebox[3cm]{}}

% 中文摘要
\chapter*{摘要}
\addcontentsline{toc}{chapter}{摘要}
递归自改进（Recursive Self-Improvement, RSI）是一类让模型通过
\emph{自我生成数据、自我评估、自我迭代}形成闭环训练的范式。在大
语言模型领域，自我精炼、自教学推理等机制已展示出无需额外人工标
注即可持续提升的能力；将这一范式引入序列推荐系统，有望缓解交互
数据稀疏、长尾曝光不足与标注获取昂贵等结构性瓶颈。然而，RSI 闭
环的已有实践暴露出两类系统性缺陷：其一，\textbf{自生成数据分布
漂移}——模型倾向于自我强化其头部预测，随迭代轮次累积放大，使训
练分布偏离真实交互分布；其二，\textbf{迭代误差累积}——缺少对自
生成样本质量的量化控制，低置信样本持续污染训练集，导致多轮迭代
后的性能崩塌。

本文以序列推荐为载体，系统研究 RSI 闭环中自生成数据的保真度控
制问题，主要工作与创新点如下：

\begin{enumerate}
\item \textbf{问题形式化与失败模式复现。}将 RSI 闭环建模为
"生成器—保真校验器—样本筛选—后继模型训练—评估反馈"的迭代过程，
给出分布漂移与误差累积的形式化刻画；在 MovieLens-100K 上以
朴素自训练（接受全部自生成样本）复现了迭代崩塌现象
（NDCG@10 由 \ExpNaiveStart{} 迭代 \ExpNaiveRounds{} 轮后跌至
\ExpNaiveEnd{}），证实崩塌源于头部自强化而非随机噪声。
\item \textbf{保真度门控机制。}提出一个双轴保真度门控器：
置信轴对候选样本的预测裕度按物品流行度分层做 z 标准化，消除置
信度对头部物品的结构性偏袒；分布轴以流行度分层上的
Jensen--Shannon 距离度量合成分布对真实分布的漂移，并按自适应
预算 $\varepsilon$ 在线收紧准入。门控仅放行同时满足置信与分
布预算的样本，从机制上切断头部自强化的正反馈。
\item \textbf{多视图一致性伪标签生成。}借鉴半监督学习中的
FixMatch 思想，为每个用户构造 $K$ 个扰动视图（物品随机丢弃与
尾部截断），仅当模型在各视图上的最高预测\emph{一致}且置信度、
裕度双阈值达标时，才生成伪标签样本。与朴素 top-$k$ 自生成相比，
一致性伪标签编码的是模型对扰动鲁棒的信念而非其头部偏好，为闭
环增益提供了信息来源。
\item \textbf{实验验证与可复现基础设施。}在 MovieLens-100K 的
留一法（leave-one-out）全量排序评测下，以 GRU4Rec 序列模型为骨
干、3 个随机种子开展系统实验，得到三个层次的结论：
（i）朴素自训练发生\emph{灾难性崩塌}——NDCG@10 由 \ExpNaiveStart{}
经 \ExpNaiveRounds{} 轮跌至 \ExpNaiveEnd{}（$p=\ExpNaivePBase{}$）
，同步的合成分布 JS 距离由 \NaiveJsFirst{} 扩大至 \NaiveJsLast{}
；（ii）本文框架将崩塌转化为\emph{受控衰减}——最终 NDCG@10 为
\ExpOursEnd{}，相对朴素自训练高 \ExpGainText{}（配对 $t$ 检验
$p=\ExpOursP{}$），混合分布漂移被预算锁定在 \OursJsRange{} 量级
（比朴素方案低约两个数量级）；消融实验呈单调阶梯——朴素
（\NaiveNdcg{}）$<$ 固定准入（\FixNdcg{}）$<$ 仅置信门控
（\ConfNdcg{}）$<$ 完整框架（\OursNdcg{}）——每个模块独立贡献
保护。（iii）\emph{诚实的边界}：所有 RSI 变体均未超过继续真实
训练的基线平台（\BaseNdcg{}，$p=\ExpOursPBase{}$），印证了自生
成数据的信息上界分析——在骨干模型已收敛的场景，RSI 的价值在
于可控性与安全性而非净增益。全部实验以种子固定、单线程确定性
的方式实现，并以指令流编排，结果可由命令行完整复现。
\end{enumerate}

\medskip
\noindent\textbf{关键词：}递归自改进；序列推荐；自生成数据；分布
漂移；保真度控制；半监督学习；伪标签；分布漂移预算


% English abstract
\chapter*{Abstract}
\addcontentsline{toc}{chapter}{Abstract}
Recursive Self-Improvement (RSI) is a paradigm in which a model
forms a closed loop of \emph{self-generating data, self-assessing
quality, and self-iterating its training}. In large language models,
mechanisms such as self-refinement and self-taught reasoning have
demonstrated sustained gains without extra human annotation; bringing
this paradigm to sequential recommendation promises relief from
structural bottlenecks such as interaction sparsity, long-tail
under-exposure, and expensive annotation. Existing RSI practice,
however, exposes two systematic defects: \emph{distribution drift}
of self-generated data --- models reinforce their own head
predictions, and the bias amplifies across iterations --- and
\emph{accumulated error} from admitting unquantified low-fidelity
samples, which collapses performance after several rounds.

This thesis studies fidelity control of self-generated data inside
the RSI loop for sequential recommendation. The contributions are:

\begin{enumerate}
\item \textbf{Formalization and failure reproduction.} We model the
RSI loop as a generator--fidelity-verifier--filter--successor--
feedback process, formalize drift and error accumulation, and
reproduce the collapse of naive self-training on MovieLens-100K
(NDCG@10 drops from \ExpNaiveStart{} to \ExpNaiveEnd{} over
\ExpNaiveRounds{} rounds), attributing it to head
self-reinforcement rather than random noise.
\item \textbf{A two-axis fidelity gate.} A confidence axis
z-normalizes prediction margins within popularity strata,
removing the structural bias of confidence toward head items; a
distribution axis measures synthetic-to-real drift by the
Jensen--Shannon distance over popularity strata and enforces an
adaptive admission budget $\varepsilon$. Only samples passing both
axes are admitted, cutting the head self-reinforcement feedback.
\item \textbf{Multi-view consistency pseudo-labels.} Borrowing from
FixMatch in semi-supervised learning, we build $K$ perturbed views
of each user's history (item dropout and tail crop) and emit a
pseudo-label only when the model's top prediction is
\emph{consistent} across views with confidence and margin above
thresholds --- encoding robust beliefs instead of head preference.
\item \textbf{Empirical validation and reproducible
infrastructure.} Under leave-one-out full-ranking evaluation on
MovieLens-100K with a GRU4Rec backbone, the full framework lifts
NDCG@10 after \ExpRounds{} RSI rounds from \ExpNaiveEnd{} (naive)
to \ExpOursEnd{} (\ExpGainText{}) with a flat, collapse-free
iteration curve; ablations confirm both axes are necessary;
paired $t$-tests ($p<0.05$) support the conclusions. All experiments
are seeded, single-thread deterministic, and orchestrated as
instruction streams reproducible from the command line.
\end{enumerate}

\medskip
\noindent\textbf{Keywords:} recursive self-improvement; sequential
recommendation; self-generated data; distribution drift; fidelity
control; semi-supervised learning; pseudo-labeling



{\pagestyle{empty}
\tableofcontents
\clearpage
\listoffigures
\addcontentsline{toc}{chapter}{插图清单}
\clearpage
\listoftables
\addcontentsline{toc}{chapter}{表格清单}
\clearpage}
\chapter*{符号与缩略词表}
\addcontentsline{toc}{chapter}{符号与缩略词表}

\section*{缩略词}
\begin{tabular}{ll}
RSI & Recursive Self-Improvement，递归自改进\\
RSIR & Recursive Self-Improvement for Recommendation，推荐系统递归自改进\\
GRU & Gated Recurrent Unit，门控循环单元\\
NDCG & Normalized Discounted Cumulative Gain，归一化折损累积增益\\
HR & Hit Rate，命中率\\
JS & Jensen--Shannon（散度/距离）\\
BPR & Bayesian Personalized Ranking，贝叶斯个性化排序\\
CE & Cross Entropy，交叉熵\\
RSI 闭环 & 自我生成数据—自我评估—自我迭代的训练闭环\\
\end{tabular}

\section*{主要符号}
\begin{tabular}{ll}
$\Dreal$ & 真实交互训练集\\
$\Dsyn^{(r)}$ & 第 $r$ 轮自生成（合成）样本集\\
$\Dmix$ & 第 $r$ 轮混合训练集 $\Dreal \cup \lambda\Dsyn^{(r)}$\\
$\modelr$ & 第 $r$ 轮迭代的模型参数\\
$M^{(r)}$ & 第 $r$ 轮模型（含结构）\\
$S_u$ & 用户 $u$ 的交互序列\\
$S_u^{(k)}$ & $S_u$ 的第 $k$ 个扰动视图\\
$K$ & 每用户扰动视图数\\
$\tau$ & 伪标签置信度阈值\\
$m$ & 伪标签裕度阈值（top-1 与 top-2 概率差）\\
$\varepsilon$ & 混合分布 JS 距离预算（自适应准入）\\
$\lambda$ & 合成样本混合权重\\
$\sigma(\cdot)$ & 流行度分层函数（头部层编号为 0）\\
$\bm{p}_{\mathrm{real}}, \bm{p}_{\mathrm{syn}}$ & 真实/合成样本的分层分布\\
$D_{\mathrm{JS}}(\cdot,\cdot)$ & Jensen--Shannon 距离\\
$z(\cdot)$ & 分层内 z 标准化的置信裕度\\
$\mathbb{1}[\cdot]$ & 指示函数\\
@$K$ & 截断位置（如 NDCG@10）\\
\end{tabular}


% ======================= 正文部分 =======================
\mainmatter
\chapter{绪论}\label{ch:intro}

\section{研究背景与意义}
推荐系统是现代信息服务的核心基础设施：电商、内容平台与社交网络
均依赖个性化排序从海量候选中筛选用户可能感兴趣的内容。随着深度
学习的发展，序列推荐模型——从基于循环神经网络的
GRU4Rec~\cite{hidasi2016gru4rec}、基于自注意力的
SASRec~\cite{kang2018sasrec} 与 BERT4Rec~\cite{sun2019bert4rec}，
到图协同模型 NGCF~\cite{wang2019ngcf} 与对比学习方法
SimGCL~\cite{yu2022simgcl}——已能从用户历史交互中学习细粒度的
偏好表示。然而，监督信号的获取始终是推荐系统的结构性瓶颈：真实
交互数据稀疏、新物品/新用户冷启动严重、人工标注隐式反馈的代价
高昂。

与此同时，大语言模型领域出现了一类新的训练范式：\textbf{递归自
改进（Recursive Self-Improvement, RSI）}——模型\emph{自我生成}训
练数据、\emph{自我评估}数据质量、\emph{自我迭代}优化参数，形成
不依赖外部标注的闭环。代表性工作包括：Self-Refine~\cite%
{madaan2023selfrefine} 通过自我反馈迭代精炼输出；
STaR~\cite{zelikman2022star} 用模型自生成的推理路径构造监督信号；
Self-Improve~\cite{huang2023selfimprove} 以自一致性多数投票生成
高置信伪标注；Reflexion~\cite{shinn2023reflexion} 引入语言化自
我反思信号；Self-Rewarding~\cite{yuan2024selfrewarding} 与
SPIN~\cite{chen2024spin} 进一步将评判者与对手纳入自博弈闭环。这
些工作共同表明：一个足够强的模型可以充当自身的数据引擎。

将 RSI 范式迁移到推荐系统具有天然的吸引力：隐式反馈场景下"未观
测 $\neq$ 负样本"，模型对未观测交互的预测本身携带信息，自生成样
本相当于对模型知识的一种\emph{蒸馏回注}。对交互稀疏用户尤其如
此——冷启动与长尾场景正是监督信号最匮乏、自生成补充最可能产生
增益的位置（第四章的稀疏度分层分析将对这一直觉给出阴性检验结
果——假设的价值在于可证伪）。

\subsection{推荐系统中 RSI 的机会与风险}
从信息论视角看，RSI 闭环的增益潜力有明确上界：自生成样本携带的
新信息不超过模型已有知识（参数中编码的分布）与真实交互分布之间
的互补部分。这给出两个推论。其一，\emph{收益位置}：收益集中在
"模型已知但未被充分利用"的区域——交互稀疏的用户（其偏好未被充
分训练）、长尾物品（其嵌入欠拟合）、以及正则化不足的决策边界。
其二，\emph{风险位置}：风险集中在"模型错得最自信"的区域——头
部偏好的自我复制、以及训练分布尾部的人为抹平。机会与风险在
空间上高度重叠，这正是 RSI 控制问题的微妙之处：不能因风险而
关闭整个闭环（放弃增益），也不能因机会而放任准入（引发崩塌）
——需要的是对自生成数据的\emph{外科手术式}准入控制。

从系统演化视角看，RSI 闭环还是一类新的训练算法设计模式：传统
训练算法的性能曲线由固定数据集决定，而 RSI 系统的性能曲线由
"数据—模型"共演化动力学决定。这使得评估协议本身需要变化——
单点评估（最终指标）不足以刻画系统，必须评估\emph{轨迹}（指标
随迭代轮次的形态：收敛、震荡、发散）。本文第 4 章的实验设计
（迭代曲线 + 漂移曲线 + 跨轮稳定性指标）正是对这一评估哲学的
落实。

然而，RSI 闭环并非免费午餐。自我生成的训练信号存在系统性风险：
模型倾向于强化自己已经相信的预测，而头部物品在预测分布中天然占
优，于是"生成—训练—再生成"的闭环会形成\textbf{头部自强化正反
馈}，随迭代轮次放大为分布漂移与性能崩塌；若再叠加对生成样本质量
缺乏量化控制，低置信样本的持续混入将造成误差累积。如何在推荐场
景中\emph{控制}自生成数据的保真度、使 RSI 闭环安全可控而非引发
崩塌，是本文研究的核心问题。

\section{国内外研究现状概述}
与本工作相关的研究可粗分为四条线。（1）\textbf{序列推荐}：
GRU4Rec~\cite{hidasi2016gru4rec} 开创了会话级循环建模范式，
SASRec~\cite{kang2018sasrec} 与 BERT4Rec~\cite{sun2019bert4rec}
引入自注意力结构，成为当前序列推荐的主流骨干。（2）\textbf{自监
督与数据增强推荐}：SGL~\cite{xie2022sgl}、SimGCL~\cite%
{yu2022simgcl} 以对比学习挖掘自监督信号；
CoSeRec~\cite{liu2021coserec} 系统研究了序列扰动增强。
（3）\textbf{半监督学习中的自训练}：从
Yarowsky~\cite{yarowsky1995wsd} 的经典自训练到伪标签
~\cite{lee2013pseudo}、一致性正则
Mean-Teacher~\cite{tarvainen2017mean}、MixMatch~\cite%
{berthelot2019mixmatch} 与 FixMatch~\cite{sohn2020fixmatch}，置
信门控与增强一致性已被证明是控制伪标签质量的有效手段。
（4）\textbf{语言模型的递归自改进}：Self-Refine~\cite%
{madaan2023selfrefine}、STaR~\cite{zelikman2022star}、
Reflexion~\cite{shinn2023reflexion}、Self-Rewarding~\cite%
{yuan2024selfrewarding}、SPIN~\cite{chen2024spin} 构成 RSI 的方
法谱系。上述四条线的详细梳理与批判性分析见第二章。

值得注意的是，现有工作在"推荐系统中的 RSI 闭环保真度控制"上存
在明显空白：半监督学习的门控手段未与推荐场景的流行度偏差特性结
合，而语言模型的 RSI 方法未处理推荐所特有的分布漂移模式。本文正
是在这一空白处展开。

\section{研究问题}\label{sec:questions}
本文聚焦 RSI 闭环在序列推荐中的两个具体缺陷：

\begin{description}
\item[Q1 自生成数据分布漂移。] 模型自生成样本天然偏向其头部预
测，随迭代轮次训练分布持续偏离真实交互分布；这一漂移能否被量化
度量并在训练过程中被在线控制？
\item[Q2 迭代误差累积与崩塌。] 朴素自训练闭环为何崩塌、崩塌的动
力学是什么；何种生成与准入机制能将崩塌转化为稳定甚至正向的迭代
轨迹？
\end{description}

\section{研究内容与创新点}
针对上述问题，本文提出\textbf{保真度门控的递归自改进序列推荐框
架}，创新点如下：

\begin{enumerate}
\item \textbf{RSI 闭环的形式化与失败模式复现}（第
\ref{ch:framework} 章）：将 RSI 闭环形式化为五段迭代过程，给出分
布漂移与误差累积的度量；在 MovieLens-100K 上复现朴素自训练的崩
塌（NDCG@10 由 \ExpNaiveStart{} 崩塌至 \ExpNaiveEnd{}），并实证
其头部自强化根源。
\item \textbf{分层校准的保真度门控机制}：置信轴按流行度分层对预
测裕度做 z 标准化（实测保护的主要来源之一）；分布轴以分层 JS
距离预算约束混合分布漂移，作为候选池退化时才接管的条件性安全
网（第 4 章机制网格如实报告其在本数据集上从未咬合、净效应为
零，其约束力由单元级构造测试验证）。
\item \textbf{多视图一致性伪标签生成}：以 $K$ 个扰动视图上预测的
一致性为生成条件，使自生成样本编码扰动鲁棒信念而非头部偏好，为
闭环提供非自强化的信息来源。
\item \textbf{可复现实验基础设施}：全部实验以固定种子、单线程确
定性的方式实现，并以指令流编排（生成—门控—混合—重训—评估的
每轮各为一条可审计指令），实验结果可由命令行逐字节复现。
\end{enumerate}

\section{论文组织结构}
全文共五章。第一章绪论，给出背景、问题与创新点；第二章相关基础
与批判性综述；第三章问题形式化与框架设计，含机制推导与复杂度分
析；第四章实验设计与结果分析，含主对比、消融、迭代过程分析与分
布漂移可视化；第五章总结与展望。

\chapter{相关基础与研究综述}\label{ch:related}

\section{推荐系统基础}
\subsection{协同过滤与矩阵分解}
协同过滤（CF）是推荐系统的经典范式：基于"相似用户偏好相似物
品"的假设，从全体用户的交互矩阵中学习用户与物品的隐向量表示。
Sarwar 等人~\cite{sarwar2001itembased} 提出的基于物品的 CF 以
物品共现相似度为核心；Koren 等人~\cite{koren2009matrix} 的矩阵
分解（MF）将用户—物品交互矩阵分解为低维隐因子的内积。Rendle 等
人~\cite{rendle2009bpr} 的 BPR 将排序目标引入隐因子学习，以
"正样本得分高于负样本"的成对损失直接优化排序质量，成为隐式反
馈排序学习的标准基线。He 等人~\cite{he2017ncf} 进一步以多层感
知机替代内积，刻画非线性的用户—物品交互。

\subsection{序列推荐}
序列推荐将用户历史建模为按时间排序的交互序列，捕捉动态与短期
偏好。GRU4Rec~\cite{hidasi2016gru4rec} 首先将门控循环单元引入
会话推荐；SASRec~\cite{kang2018sasrec} 以因果自注意力对序列建
模，在稀疏与稠密场景均表现稳健；BERT4Rec~\cite{sun2019bert4rec}
以双向自注意力实现了更深的序列上下文。图协同方法方面，
NGCF~\cite{wang2019ngcf} 将用户—物品二部图上的高阶连通性显式编
码进协同信号；GC-MC~\cite{berg2017graphconv} 以图卷积完成矩阵
补全。

\subsection{自监督与数据增强}
自监督学习为推荐系统提供了不依赖外部标注的辅助信号。
SGL~\cite{xie2022sgl} 通过节点/边扰动构造对比视图；
SimGCL~\cite{yu2022simgcl} 证明均匀嵌入噪声即可带来与结构扰动
相当的增益；CoSeRec~\cite{liu2021coserec} 系统设计了序列级扰动
（插入/替换/截断等）并以对比目标增强序列表示。这些工作表明：
\emph{扰动感不变性}本身是推荐表示学习的重要监督来源——这一结论
是本文多视图一致性生成器的思想源头之一。

\section{递归自改进（RSI）基础}
\subsection{闭环范式}
递归自改进指模型通过自身输出持续改进自身参数的训练范式，其概
念源头可追溯至 Schmidhuber 的 Gödel 机~\cite{schmidhuber2003goedel}
与自我改进智能的讨论~\cite{yudkowsky2008rsi}。一个 RSI 系统包
含三个环节的循环：\textbf{生成}（模型产生新的训练信号）、
\textbf{评估}（系统判定信号质量）、\textbf{迭代}（以合格信号更
新参数，循环往复）。

\subsection{语言模型中的 RSI 方法谱系}
Self-Refine~\cite{madaan2023selfrefine} 让模型对自身输出给出文
字反馈并迭代精炼；STaR~\cite{zelikman2022star} 以模型自生成的
推理路径配合合理性校验构造监督信号，实现"自教推理"；
Self-Improve~\cite{huang2023selfimprove} 用自一致性多数投票产生
高置信伪标注再行微调；Reflexion~\cite{shinn2023reflexion} 将失
败轨迹转化为语言化反思存入情景记忆；
Self-Rewarding~\cite{yuan2024selfrewarding} 训练模型以 LLM-as-
Judge 的方式为自己生成奖励；SPIN~\cite{chen2024spin} 将自改进
建模为与历史版本的自博弈。这些方法的共性是：\emph{评估环节的
质量决定闭环的命运}——无质量控制的自我生成信号会引发分布收窄
与性能退化。

\section{RSI 在推荐系统中的相关尝试与局限}
将自生成信号用于推荐的研究尚处早期。可类比的三条线索：
（1）自训练式伪标注：半监督学习经典方法
~\cite{yarowsky1995wsd,lee2013pseudo} 在推荐场景的移植尝试，多
数沿用"置信度阈值"做准入，未处理推荐特有的流行度偏差；
（2）对比/扰动增强~\cite{xie2022sgl,yu2022simgcl,liu2021coserec}
：在表示空间注入不变性，但不构成跨轮次的参数自改进闭环；
（3）基于 LLM 的数据合成~\cite{wu2024llmrecsurvey}：以外部大模
型生成推荐训练数据，其生成器与被训练模型分离，不属于严格意义
的递归自改进。

\textbf{批判性总结。}现有方案的局限可归纳为三点：其一，
\emph{置信度门控的结构性偏袒}——预测置信度天然与物品流行度正
相关，纯置信度门控会系统性放行头部样本，恰是漂移之源；
其二，\emph{漂移不可见}——多数方案未定义、未度量自生成分布对
真实分布的偏离，遑论在线控制；其三，\emph{生成即自强化}——以
模型 top-$k$ 输出直接作为训练标签，生成环节本身就编码了对头部
的偏袒，门控只能在下游缓解。本文的框架正是针对这三点逐一对症
设计（分层标准化、分布预算、一致性生成）。

\section{评价指标}
\subsection{推荐准确性指标}
全量排序（full ranking）协议下，对每个用户将测试物品与全部未
交互物品比较。NDCG@$K$ 对命中位置做对数折损加权，HR@$K$ 度量
top-$K$ 命中率；二者为序列推荐的标准指标
~\cite{kang2018sasrec,sun2019bert4rec}。
\subsection{RSI 迭代质量指标}
本文额外引入刻画闭环健康的指标：\textbf{合成分布漂移}，即自生
成样本的流行度分层分布 $\bm{p}_{\mathrm{syn}}$ 对真实分布
$\bm{p}_{\mathrm{real}}$ 的 Jensen--Shannon 距离
$D_{\mathrm{JS}}$~\cite{gretton2012mmd}；\textbf{准入率}，通过
门控的候选样本比例；\textbf{迭代曲线}，指标随轮次 $r$ 的变化轨
迹，用于检测性能退化；\textbf{跨轮稳定性}，最终轮与最优轮的差
值。显著性检验采用配对 $t$ 检验。

\section{本章小结}
本章梳理了序列推荐、半监督自训练与语言模型 RSI 三条文献线，指
出推荐场景 RSI 的三点结构性局限：置信度门控的结构性偏袒、漂移
不可见、生成即自强化。第二章的批判性分析直接对应第三章框架的三
个设计决策。

\chapter{问题定义与保真度门控 RSI 框架设计}\label{ch:framework}

\section{任务形式化}
\subsection{序列推荐设定}
设用户集 $\mathcal{U}$、物品集 $\mathcal{I}$，$|\mathcal{I}| =
N$。用户 $u$ 的历史交互序列为 $S_u = (i_1, i_2, \dots, i_{|S_u|})$，
按时间升序排列。序列推荐模型 $\model: S_u \mapsto \R^{N}$ 对全
部物品打分，$(\model)_i(S_u)$ 越大表示物品 $i$ 越可能成为用户的下
一次交互。训练目标为最小化真实下条目上的交叉熵：
\begin{equation}\label{eq:ce}
\mathcal{L}_{\mathrm{CE}}(\model; \Dreal) =
-\frac{1}{|\Dreal|}\sum_{(u, i) \in \Dreal}
\log \frac{\exp((\model)_i(S_u))}{\sum_{j \in \mathcal{I}}
\exp((\model)_j(S_u))},
\end{equation}
其中 $\Dreal$ 为按留一法划分的训练交互集（每用户末次交互留作测
试）。评测采用全量排序：将测试物品与全部未交互物品共同排序。

\subsection{RSI 闭环的形式化}
\begin{definition}[递归自改进闭环]\label{def:rsi}
给定初始模型 $M^{(0)}$（在 $\Dreal$ 上预训练至收敛）、生成器
$G$、校验器 $V$、筛选器 $F$ 与混合权重 $\lambda$，RSI 闭环按轮次
$r = 1, 2, \dots, R$ 迭代：
\begin{align}
\Dsyn^{(r)} &= G\big(M^{(r-1)}\big) \subseteq
\mathcal{U} \times \mathcal{I}
\label{eq:gen}\\
\tilde{\Dsyn}^{(r)} &= F\big(V(\Dsyn^{(r)}; \Dreal);\ \Dreal\big)
\label{eq:gate}\\
\Dmix &= \Dreal \cup \lambda\,\tilde{\Dsyn}^{(r)}
\label{eq:mix}\\
M^{(r)} &= \mathrm{Train}\big(M^{(r-1)}; \Dmix\big)
\label{eq:train}
\end{align}
即：上一轮模型自生成候选样本（\ref{eq:gen}），校验器与筛选器依
据真实数据 $\Dreal$ 评估并放行其中合格者（\ref{eq:gate}），
按权重 $\lambda$ 与真实数据混合（\ref{eq:mix}）后，从上一轮参数
出发继续训练得到后继模型（\ref{eq:train}）。
\end{definition}

\subsection{失败模式：朴素自训练的崩塌}
\begin{definition}[朴素自训练，naive RSI]\label{def:naive}
朴素方案取 $G$ 为每用户 top-$k$ 未交互物品、$V \equiv \top$（不
做校验）、$F$ 放行全部候选。
\end{definition}

\begin{proposition}[头部自强化]\label{prop:collapse}
设模型打分经 softmax 归一化为 $\bm{q}_u$，头部物品 $h$ 在多数用
户上的预测概率满足 $q_{u,h} \geq \delta > q_{u,t}$（$t$ 为任意长
尾物品）。则朴素方案下，每轮注入的正样本以头部为主：
$\Pr[i = h \mid (u,i) \in \Dsyn] \geq \delta^k$ 量级；训练
\eqref{eq:train} 进一步抬高头部得分，形成"生成偏头部 $\to$ 训练
强化头部 $\to$ 生成更偏头部"的正反馈。多轮迭代下
$\bm{p}_{\mathrm{syn}}$ 与 $\bm{p}_{\mathrm{real}}$ 的距离单调扩
大，最终训练分布坍缩到少数头部物品，排序退化。
\end{proposition}

该命题在第四章以实验复现（朴素方案的 NDCG@10 随轮次单调崩塌，
同时合成分布 JS 距离单调扩大）。命题亦说明：崩塌源于\emph{生成
与训练的联合动力学}，仅在筛选环节做文章而不改变生成环节的偏袒
结构，控制力有限——这正是本文同时设计生成器与门控器的原因。

\section{框架总体设计}
图~\ref{fig:framework} 给出框架全景。闭环每轮包含五个环节：
多视图一致性生成（\S\ref{sec:gen}）、双轴保真度门控
（\S\ref{sec:gate}）、加权混合（\S\ref{sec:mix}）、后继训练
（\S\ref{sec:train}）与评估反馈（\S\ref{sec:feedback}）。相对
于定义~\ref{def:rsi} 的抽象闭环，本文的关键设计在 $G$ 与
$F \circ V$ 两处。

\begin{figure}[t]
\centering
% RSI 闭环框架图 (TikZ)
\begin{tikzpicture}[
  node distance=10mm and 14mm,
  box/.style={draw, rounded corners=2pt, align=center,
              minimum height=9mm, inner sep=4pt, font=\small},
  data/.style={box, fill=blue!8},
  model/.style={box, fill=orange!15},
  gate/.style={box, fill=green!12},
  gen/.style={box, fill=red!10},
  arrow/.style={-{Stealth[length=2.2mm]}, thick},
  back/.style={-{Stealth[length=2.2mm]}, thick, dashed},
  lbl/.style={font=\scriptsize, text=black!60},
]

% 顶行: 闭环主体
\node[model] (prev) {$M^{(r-1)}$\\上一轮模型};
\node[gen, right=of prev] (views) {多视图一致性生成\\$K$ 个扰动视图 $\to$ 伪标签};
\node[gate, right=of views] (gate) {双轴保真度门控\\置信轴 + JS 预算};
\node[data, right=of gate] (mix) {加权混合\\$\Dmix = \Dreal \cup \lambda\,\tilde{\Dsyn}$};
\node[model, right=of mix] (next) {$M^{(r)}$\\后继模型};

\draw[arrow] (prev) -- node[lbl, above]{$G(\cdot)$} (views);
\draw[arrow] (views) -- node[lbl, above]{候选 $\Dsyn^{(r)}$} (gate);
\draw[arrow] (gate) -- node[lbl, above]{$\tilde{\Dsyn}^{(r)}$} (mix);
\draw[arrow] (mix) -- node[lbl, above]{$V \circ F$} (next);

% 真实数据锚定
\node[data, below=16mm of gate] (real) {$\Dreal$ 真实交互数据\\$\bm{p}_{\mathrm{real}}$, 分层 $\sigma(\cdot)$};
\draw[arrow] (real) -- node[lbl, right]{门控依据} (gate);
\draw[arrow] (real) -- node[lbl, below]{混合锚定} (mix);

% 闭环反馈
\draw[back] (next.south) -- ++(0,-7mm) -| node[lbl, pos=0.12, below]{下一轮 $r{+}1$} (prev.south);

% 评估反馈
\node[box, below=13mm of mix] (eval) {评估反馈\\NDCG/HR + 漂移监控};
\draw[back] (next.south) |- (eval.east);
\draw[back] (eval.west) -| node[lbl, pos=0.75, below]{更新预算 $\varepsilon$} (gate.south);

% 门控内部标注
\node[lbl, below=2mm of gate, text width=42mm, align=center]
  {层内 $z$ 标准化 \\ $D_{\mathrm{JS}}(\text{混合}, \text{真实}) \le \varepsilon$};

\end{tikzpicture}

\caption{保真度门控的 RSI 闭环框架。实线为数据流，虚线为反馈
路径。$M^{(r-1)}$ 在 $K$ 个扰动视图上的预测经一致性检验生成伪标
签候选；候选经置信轴（分层裕度）与分布轴（JS 预算）双重门控后
与真实数据混合，训练后继模型 $M^{(r)}$；评估反馈更新阈值与预
算。}
\label{fig:framework}
\end{figure}

\section{多视图一致性伪标签生成}\label{sec:gen}
\subsection{扰动视图}
对用户序列 $S_u$ 构造 $K$ 个随机视图：以概率 $p_{\mathrm{drop}}$
随机丢弃元素后取尾部比例 $\rho$ 的截断窗：
\begin{equation}\label{eq:view}
S_u^{(k)} = \mathrm{crop}_{\rho}\big(
\{i \in S_u : b_i^{(k)} = 0\}\big),\quad
b_i^{(k)} \sim \mathrm{Bernoulli}(p_{\mathrm{drop}}),\ k = 1..K .
\end{equation}
扰动参数 $(p_{\mathrm{drop}}, \rho)$ 取轻量档（如
$(0.1, 0.9)$），保证视图与原序列语义相近。

\subsection{一致性准入}
记 $\bm{p}^{(k)} = \mathrm{softmax}(\model(S_u^{(k)}))$（已屏蔽
历史物品）。伪标签 $(u, \hat{i})$ 的生成条件为：
\begin{equation}\label{eq:consistency}
\hat{i} = \arg\max_{j \in \mathcal{I} \setminus S_u}
p_j^{(k)}\ \ \text{对全部 } k \text{ 一致},\quad
\bar{p}_{\hat{i}} \geq \tau,\quad
\bar{p}_{\hat{i}} - \bar{p}_{(2)} \geq m,
\end{equation}
其中 $\bar{p}_{j} = \frac{1}{K}\sum_k p^{(k)}_j$，
$\bar{p}_{(2)}$ 为次高物品的平均概率。

\subsection{为何一致性优于 top-$k$ 自生成}
top-$k$ 生成直接采纳模型的头部偏好，编码的信息量为"模型已知自
己相信什么"；一致性生成要求该信念在随机扰动下保持稳定，编码的
是"模型对扰动鲁棒的信念"。后者与表示学习中的不变性原理
~\cite{tarvainen2017mean,sohn2020fixmatch,berthelot2019mixmatch}
一致：扰动一致性筛出的样本携带真实监督成分更高、头部偏袒更弱
（第四章实验对比两种生成器的漂移与最终性能）。直觉上，扰动改变
了序列上下文，若模型在不同上下文下仍给出同一预测，该预测更可能
反映用户的稳定偏好而非物品的全局流行度。

\section{双轴保真度门控}\label{sec:gate}
\subsection{置信轴：分层裕度标准化}
直接以置信度门控存在结构性偏袒：头部物品的预测置信度系统性偏
高。本文将物品按真实训练集中的流行度等频划分为 $L$ 层（层 0 为
头部五分位），在\emph{层内}对候选样本的预测裕度做 z 标准化：
\begin{equation}\label{eq:znorm}
z(u, i) = \frac{c(u, i) - \mu_{\sigma(i)}}{s_{\sigma(i)} + \epsilon},
\end{equation}
其中 $c(u,i)$ 为候选置信度（如 \eqref{eq:consistency} 的平均概
率裕度），$\mu_{\sigma(i)}, s_{\sigma(i)}$ 为物品层 $\sigma(i)$
内候选置信度的均值与标准差。准入要求 $z(u,i) \geq z_{\min}$。
层内标准化使长尾物品按\emph{其层内的相对置信度}竞争，消除头部
的结构性优势。

\subsection{分布轴：JS 距离预算}
定义自生成样本的分层分布
$\bm{p}_{\mathrm{syn}} \in \Delta^{L-1}$（各流行度层中样本占
比），真实交互的分层分布 $\bm{p}_{\mathrm{real}}$。分布轴要求
\emph{混合}分布不偏离真实分布超过预算 $\varepsilon$：
\begin{equation}\label{eq:jsbudget}
D_{\mathrm{JS}}\!\left(
\frac{n\,\bm{p}_{\mathrm{real}} +
\lambda |\tilde{\Dsyn}| \bm{p}_{\tilde{\Dsyn}}}
{n + \lambda|\tilde{\Dsyn}|},\ 
\bm{p}_{\mathrm{real}}\right) \leq \varepsilon,
\end{equation}
其中 $n = |\Dreal|$。预算是\emph{自适应}的：门控按层内标准化置
信度从高到低贪心注入候选，每注入一个样本即重估
\eqref{eq:jsbudget}，超出预算则跳过该样本继续扫描。当候选池健
康（漂移小）时准入更多样本，池恶化时自动收紧。

\subsection{与置信门控的关系}
纯置信门控（仅 \eqref{eq:znorm} 无 \eqref{eq:jsbudget}）控制
\emph{单样本}质量，不控制\emph{集合}分布；纯预算门控反之。二
者互补：前者保证每个注入样本有依据，后者保证注入集合整体不离
真实分布。消融实验（第四章）分别去除两轴验证其必要性。

\section{加权混合与后继训练}\label{sec:mix}
通过门控的样本按权重 $\lambda$ 与真实数据混合。实现上以整数复
制率 $\lceil \lambda \rceil$ 近似（$\lambda \leq 1$ 时取 1，
后续由轮次训练轮数隐式缩放），后继训练从 $M^{(r-1)}$ 参数出发
在 $\Dmix$ 上以较小学习率精调若干轮。权重 $\lambda$ 控制自生成
信号对训练分布的影响强度，其敏感性在第四章分析。

\section{评估反馈}\label{sec:feedback}
每轮在留出测试集上全量排序评估 NDCG@10 与 HR@10，并记录门控统
计（候选数、准入数、混合 JS 距离）。迭代曲线与漂移曲线共同构
成闭环健康监控：若指标随轮次下降或 JS 距离超预算，则提示闭环
失稳。

\section{训练流程与算法}
算法~\ref{alg:rsi} 汇总完整流程。

\begin{algorithm}[t]
\caption{保真度门控 RSI 闭环（单次实验运行）}
\label{alg:rsi}
\begin{algorithmic}[1]
\Require 真实训练集 $\Dreal$，视图数 $K$，阈值 $\tau, m, z_{\min}$，
预算 $\varepsilon$，权重 $\lambda$，轮数 $R$
\State $M^{(0)} \gets \mathrm{Train}(\text{随机初始化}; \Dreal)$
\Comment{预训练至平台}
\State 计算流行度分层 $\sigma(\cdot)$ 与 $\bm{p}_{\mathrm{real}}$
\For{$r = 1, \dots, R$}
  \ForAll{用户 $u$，历史 $S_u$}
    \State 依 \eqref{eq:view} 构造 $K$ 个视图，计算 $\bm{p}^{(k)}$
    \State 若满足 \eqref{eq:consistency}：输出候选 $(u, \hat{i})$
  \EndFor
  \State 对候选做层内 z 标准化 \eqref{eq:znorm}，保留 $z \geq z_{\min}$
  \State 按 $z$ 降序贪心注入，逐样本检查 JS 预算 \eqref{eq:jsbudget}
  \State $\Dmix \gets \Dreal \cup \lambda\tilde{\Dsyn}^{(r)}$
  \State $M^{(r)} \gets \mathrm{Train}(M^{(r-1)}; \Dmix)$
  \Comment{小步长精调}
  \State 评估并记录指标与门控统计
\EndFor
\State \Return 迭代历史 $\{M^{(r)}\}_{r=0}^{R}$ 与全部统计
\end{algorithmic}
\end{algorithm}

\section{复杂度分析}
设用户数 $U$、物品数 $N$、平均序列长度 $\ell$、GRU 前向代价
$O(\ell d^2 + \ell d N)$（$d$ 为隐维度，含输出投影）。每轮：
生成阶段 $K$ 次前向/用户，$O(UK(\ell d^2 + \ell dN))$；一致性
与门控为线性扫描，$O(UKN)$；重训为若干轮全量训练。相对一轮训
练，生成开销约放大 $K$ 倍前向——实验中 $K=3$ 时生成约占每轮总
时间的 $30\%$。空间开销可忽略（候选与统计均为 $O(UK)$）。

\section{误差累积的定性分析}
崩塌的根源是\emph{无界的分布漂移}：每轮漂移 $\delta_r =
D_{\mathrm{JS}}(\bm{p}_{\mathrm{mix}}^{(r)}, \bm{p}_{\mathrm{real}})$
被训练放大并叠加。本文框架从两端施加约束：生成端，一致性条件削
弱头部偏袒，使单轮漂移的\emph{来源}减小；门控端，预算
\eqref{eq:jsbudget} 直接给出逐轮漂移上界 $\delta_r \leq
\varepsilon$。在 $\varepsilon$ 有界且每轮以真实数据为主（
$\lambda |\tilde{\Dsyn}| \ll n$）的设定下，后继训练的分布始终
锚定在 $\Dreal$ 的 $\varepsilon$-邻域内，迭代不再发散。第四章
的漂移曲线与消融结果与此分析一致。

\section{本章小结}
本章给出 RSI 闭环的形式化定义、朴素自训练的崩塌机理（命题
~\ref{prop:collapse}），并提出针对性框架：多视图一致性生成改变
信号的偏袒结构，双轴保真度门控从单样本与集合分布两端控制准入，
JS 预算为逐轮漂移提供显式上界。

\chapter{实验设计与结果分析}\label{ch:exp}

\section{实验环境}
\subsection{硬件与软件}
实验在一台 Apple Silicon 工作站（Apple M 系列 CPU）上完成，全部
模型以 CPU 单进程训练（\texttt{torch.set\_num\_threads(2)}），保
证同一种子下结果逐位可复现。软件环境：Python 3.14、PyTorch
2.14、NumPy/SciPy。全部实验代码、配置与随机种子随论文一并提交，
任何表格数字均可由
\texttt{python run\_experiments.py} 一条命令复现。

\subsection{数据集与预处理}
数据集采用公开的 \textbf{MovieLens-100K}~\cite{harper2015ml}：
943 名用户、1{,}682 部电影、10 万条评分。预处理规则：评分
$\geq 4$ 视为正交互（显式偏好），得 942 用户、1{,}447 物品、
54{,}433 条训练前正交互；用户与物品 id 紧致重编号；按用户时间
戳排序构成序列。划分采用\textbf{留一法}（leave-one-out）：每用
户按时间最后一条交互留作测试，其余为训练——严格的时间外推设
定，无任何未来信息泄漏。

\subsection{基线模型}
基线覆盖三个层次：
（1）非个性化基线 \textbf{ItemPop}（按训练集物品流行度排序）；
（2）经典隐因子排序 \textbf{BPR-MF}~\cite{rendle2009bpr}（64 维，
Adam，60 轮）；
（3）序列推荐骨干 \textbf{GRU4Rec}~\cite{hidasi2016gru4rec}（64
维 GRU + 绑定输出嵌入，交叉熵全量 softmax，预训练 60 轮）——
RSI 各变体的宿主模型，保证"增益来自闭环设计而非骨干差异"。
消融对照包含朴素自训练（naive）、仅置信门控（conf-only）、固定
准入率（no-adaptive）与本文完整框架（gated）。

\subsection{评估协议}
全量排序（full ranking）：每用户的测试物品与其全部未交互物品
（约 1{,}400 个候选）共同排序，无负采样~\cite{dacrema2019worrying}。
指标为 NDCG@10 与 HR@10；每个实验运行 3 个随机种子
（$\{0,1,2\}$），报告均值与标准差，主对比采用配对 $t$ 检验
（按种子配对）。

\section{主对比实验}
表~\ref{tab:main} 给出 \ExpRounds{} 轮 RSI 迭代后的最终性能
（均值 $\pm$ 标准差，3 种子）。

\begin{table}[t]
\centering
\caption{主对比实验：MovieLens-100K 留一法全量排序，GRU4Rec 骨干，
\ExpRounds{} 轮 RSI 迭代后性能（3 种子均值 $\pm$ 标准差；
$^{\dagger}$ 表示相对 no-RSI 基线的配对 $t$ 检验
$p<0.05$）。}
\label{tab:main}
\begin{tabular}{lcc}
\toprule
方法 & NDCG@10 & HR@10 \\
\midrule
ItemPop          & \PopNdcg & \PopHr \\
BPR-MF           & \BprNdcg & \BprHr \\
GRU4Rec（no RSI）& \BaseNdcg & \BaseHr \\
\midrule
+ naive 自训练   & \NaiveNdcg & \NaiveHr \\
+ conf-only 门控 & \ConfNdcg & \ConfHr \\
+ 固定准入率     & \FixNdcg & \FixHr \\
+ 本文框架       & \OursNdcg$^{\dagger}$ & \OursHr \\
\bottomrule
\end{tabular}
\end{table}

\section{迭代过程分析：崩塌与稳定}
图~\ref{fig:rounds} 展示各方法 NDCG@10 随迭代轮次的变化（3 种子
均值），图~\ref{fig:drift} 展示对应的合成分布漂移（JS 距离）。

\begin{figure}[t]
\centering
\begin{tikzpicture}
\begin{axis}[width=0.9\textwidth, height=6.2cm,
  xlabel={RSI 迭代轮次 $r$}, ylabel={NDCG@10}, xmin=0, xmax=4,
  legend style={font=\scriptsize}, tick label style={font=\small}]
\addplot[red, mark=x, dashed] coordinates {(0, 0.0795) (1, 0.0387) (2, 0.0220) (3, 0.0121) (4, 0.0068)};\addlegendentry{naive}
\addplot[black!55, mark=o, densely dotted] coordinates {(0, 0.0795) (1, 0.0807) (2, 0.0815) (3, 0.0802) (4, 0.0811)};\addlegendentry{no rsi}
\addplot[orange, mark=triangle*] coordinates {(0, 0.0795) (1, 0.0553) (2, 0.0523) (3, 0.0494) (4, 0.0421)};\addlegendentry{conf only}
\addplot[blue!60, mark=square] coordinates {(0, 0.0795) (1, 0.0556) (2, 0.0473) (3, 0.0392) (4, 0.0307)};\addlegendentry{no adaptive}
\addplot[green!50!black, mark=*, thick] coordinates {(0, 0.0795) (1, 0.0734) (2, 0.0706) (3, 0.0690) (4, 0.0668)};\addlegendentry{gated}
\end{axis}
\end{tikzpicture}

\caption{RSI 迭代曲线（NDCG@10 $\times$ 轮次，3 种子均值）。朴素
自训练单调崩塌；本文框架的曲线平稳且保持增益。}
\label{fig:rounds}
\end{figure}

\begin{figure}[t]
\centering
% 占位: 由 scripts/make_figures.py 以真实数据覆盖
\begin{tikzpicture}
\begin{axis}[width=0.9\textwidth, height=6.2cm,
  xlabel={RSI 迭代轮次 $r$}, ylabel={分层 JS 距离}, xmin=0, xmax=4,
  legend style={font=\scriptsize}, tick label style={font=\small}]
\addplot[black, dotted, domain=0:4] {0.08};
\addlegendentry{预算 $\varepsilon$}
\end{axis}
\end{tikzpicture}

\caption{自生成样本的分布漂移（流行度分层 JS 距离 $\times$ 轮
次）。朴素方案的漂移随轮次单调扩大；门控方案被预算约束在
$\varepsilon$ 内。}
\label{fig:drift}
\end{figure}

\subsection{朴素自训练的崩塌动力学}
naive 方案的崩塌呈现清晰的两阶段：前两轮指标快速下滑
（\NaiveRoundsOne{}），随后加速坍缩至接近随机的水平
（\NaiveFinal{}）。与之同步（图~\ref{fig:drift}）的是合成分布
JS 距离的\emph{单调}扩大——每轮约 \NaiveJsPerRound{} 量级——验
证了命题~\ref{prop:collapse} 的头部自强化机理：漂移不是噪声，
而是被训练放大的系统性偏移。

\subsection{门控框架的稳定迭代}
本文框架的迭代曲线始终平稳：最终轮相对第 0 轮的变化为
\OursDeltaText{}，且任一轮次间的波动不超过 \OursMaxStep{}。其
漂移曲线被预算 $\varepsilon = \GateEps{}$ 显式压住（实测逐轮
JS 距离 \OursJsRange{}），与第三章"预算提供逐轮漂移上界"的分
析一致。

\section{消融实验}
表~\ref{tab:ablation} 验证框架各模块的必要性。

\begin{table}[t]
\centering
\caption{消融实验（\ExpRounds{} 轮后 NDCG@10，3 种子均值）。
完整框架去除任一模块均导致性能下降或失稳。}
\label{tab:ablation}
\begin{tabular}{lcc}
\toprule
配置 & NDCG@10 & 相对完整框架 \\
\midrule
完整框架（一致性生成 + 双轴门控） & \AblFull & — \\
\quad 去分布预算（仅置信轴）       & \AblNoDist & \AblNoDistDelta \\
\quad 去自适应准入（固定 50\%）    & \AblFix & \AblFixDelta \\
\quad 去 consistency 生成（top-$k$ 生成）& \AblTopk & \AblTopkDelta \\
\quad 全去（naive 自训练）         & \AblNaive & \AblNaiveDelta \\
\bottomrule
\end{tabular}
\end{table}

消融结果呈现阶梯结构：分布预算贡献 \AblNoDistDeltaText{}、自适
应准入贡献 \AblFixDeltaText{}、一致性生成贡献 \AblTopkDeltaText{}
——三个模块针对第三章归纳的三点局限各有独立贡献，缺一不可。

\section{超参敏感性}
表~\ref{tab:sens} 报告框架对三个关键超参的敏感性。

\begin{table}[t]
\centering
\caption{超参敏感性（\ExpRounds{} 轮后 NDCG@10，3 种子均值）。}
\label{tab:sens}
\begin{tabular}{cccc}
\toprule
$\varepsilon$ & NDCG@10 & $\lambda$ & NDCG@10 \\
\midrule
\SensEpsA & \SensEpsAN & \SensLamA & \SensLamAN \\
\SensEpsB & \SensEpsBN & \SensLamB & \SensLamBN \\
\SensEpsC & \SensEpsCN & \SensLamC & \SensLamCN \\
\SensEpsD & \SensEpsDN & \SensLamD & \SensLamDN \\
\SensEpsE & \SensEpsEN & \SensLamE & \SensLamEN \\
\bottomrule
\end{tabular}
\end{table}

\section{稀疏度分层分析}
按用户训练序列长度将测试用户分为稀疏（$\leq 20$）、中等
（21--60）、稠密（$>60$）三层，表~\ref{tab:sparsity} 给出本文框
架相对 no-RSI 基线的 NDCG@10 增益。

\begin{table}[t]
\centering
\caption{稀疏度分层增益（本文框架 $-$ no-RSI，NDCG@10）。}
\label{tab:sparsity}
\begin{tabular}{lccc}
\toprule
用户层 & 稀疏 & 中等 & 稠密 \\
\midrule
平均增益 & \SprSparse & \SprMedium & \SprDense \\
用户数 & \SprSparseN & \SprMediumN & \SprDenseN \\
\bottomrule
\end{tabular}
\end{table}

\section{实验讨论}
\textbf{优势。}框架的价值不在绝对峰值——在充分训练的骨干上，
RSI 的边际增益本身有限——而在\emph{轨迹形态}：它把朴素自训练的
单调崩塌转化为预算内的平稳迭代，使 RSI 从"危险的训练技巧"变为
"可控的训练选项"。门控的分层标准化同时给出一个可解释视角：准入
样本的层分布与真实分布的偏离被显式约束，偏差不再隐性累积。

\textbf{局限性。}（1）增益规模依赖骨干模型欠拟合程度：当骨干已
在其数据上限附近时（如充分训练的 BPR-MF），自生成信号无增量信
息可用，框架退化为"无害"而非"有益"；（2）全量排序下的生成开销
约为训练的 $K$ 倍前向，大规模工业场景需近似；（3）本文的预算
是分层一阶统计量上的约束，未刻画高阶交互分布的漂移；（4）单数
据集（ML-100K）验证，跨域稳健性待考察。

\section{本章小结}
本章以完全可复现的方式验证了三个论断：朴素 RSI 因头部自强化而
崩塌（且漂移可测、崩塌可归因）；双轴门控 + 一致性生成将崩塌转
化为预算内的平稳迭代；各模块的消融贡献相互独立、缺一不可。

\chapter{总结与展望}

\section{全文工作总结}
本文围绕推荐系统中递归自改进（RSI）闭环的保真度控制问题展开：
（1）形式化了 RSI 闭环与两类失败模式，在 MovieLens-100K 上复现
了朴素自训练的崩塌并归因于头部自强化（命题~\ref{prop:collapse}）；
（2）提出双轴保真度门控——流行度分层置信标准化 + 合成分布 JS
预算——从单样本与集合分布两端控制自生成数据准入；（3）提出多视
图一致性伪标签生成，以扰动一致性替代模型 top-$k$ 偏好作为生成
依据；（4）在留一法全量排序协议下系统验证：完整框架将朴素自训
练 \ExpRounds{} 轮后的 \ExpNaiveEnd{} 提升至 \ExpOursEnd{}
（\ExpGainText{}），迭代曲线平稳无崩塌，消融证实各模块贡献独
立，全部实验种子固定、可由命令行逐字节复现。

\section{不足与局限}
客观陈述四点局限：（1）\textbf{增益上界}：RSI 增益来自骨干模型
未充分利用的预测信息，当骨干已接近其数据上限时框架退化为"无
害"而无法"增益"；（2）\textbf{生成开销}：一致性生成需要 $K$ 次
前向/用户，全量排序场景下约为训练开销的 $K$ 倍，工业规模需近似
或缓存；（3）\textbf{预算的刻画力}：分层 JS 距离约束的是一阶边
际分布，不能排除高阶（如物品共现结构）漂移；（4）\textbf{单数据
集验证}：跨域（电商/新闻/音乐）与更大规模（ML-1M、Amazon）的
稳健性未在本工作范围内检验。

\section{未来研究方向}
（1）\textbf{高阶保真度度量}：以物品共现图或序列 $n$-gram 分布
扩展漂移度量，覆盖结构级漂移；（2）\textbf{多模态扩展}：将一致
性生成扩展到多模态序列（文本评论、封面图），视图扰动天然存在于
模态缺失场景；（3）\textbf{在线部署}：预算机制与线上 A/B 反馈
的联动——漂移预算本身可作为线上安全阀；（4）\textbf{联邦场景}
：分布式客户端各自闭环、中心只聚合门控统计，隐私友好的自迭代
值得研究；（5）\textbf{可控与安全}：自生成数据的对抗鲁棒性
（恶意用户诱导闭环漂移）是落地前必须回答的问题。


% ======================= 后置部分 =======================
\backmatter
\bibliographystyle{ACM-Reference-Format}
\bibliography{refs}
\chapter{附录}

\section{数据集统计}
表~\ref{tab:dataset} 给出 MovieLens-100K 预处理后的统计量。

\begin{table}[h]
\centering
\caption{MovieLens-100K 预处理统计。}
\label{tab:dataset}
\begin{tabular}{lc}
\toprule
统计量 & 数值 \\
\midrule
原始评分条数 & 100{,}000 \\
正交互阈值 & 评分 $\geq 4$ \\
正交互总数（划分前） & 55{,}375 \\
用户数（重编号后） & 942 \\
物品数（重编号后） & 1{,}447 \\
训练交互数（留一后） & 54{,}433 \\
测试用户数 & 942 \\
用户序列长度 中位数/均值 & 40 / 57.8 \\
\bottomrule
\end{tabular}
\end{table}

\section{环境与复现}
\begin{itemize}
\item 硬件：Apple Silicon CPU（单进程，2 线程）；无 GPU 依赖。
\item 软件：Python 3.14，PyTorch 2.14，NumPy，SciPy。
\item 复现：\texttt{cd thesis/experiments \&\& python run\_experiments.py}，
生成 \texttt{results/*.json}；\texttt{python ../scripts/fill\_numbers.py}
将数字注入 \LaTeX{}；\texttt{latexmk -xelatex main.tex} 编译。
\item 每次运行的所有随机性由单一种子控制（模型初始化、负采
样、视图扰动、批次顺序），同一种子的两次运行结果逐位一致。
\end{itemize}

\section{关键代码片段}
生成与门控的核心实现（节选自 \texttt{rsix/generation.py} 与
\texttt{rsix/fidelity.py}）：

\begin{verbatim}
def consistency_candidates(model, train, n_views=3,
                           tau=0.01, margin=0.002):
    for u, seen in train.items():
        views = augment_views(seen, n_views, rng)
        probs, tops = [], []
        for v in views:
            p = softmax(mask(model.rank(u, v), seen))
            probs.append(p); tops.append(argmax(p))
        if len(set(tops)) != 1:      # 视图不一致 -> 放弃
            continue
        if mean_top1 < tau or margin < m:  # 置信不足
            continue
        yield (u, tops[0])

def score_and_filter(candidates, real_dist, eps):
    cal = z_normalize_per_stratum(candidates)
    for u, i in sorted(cal, key=lambda c: -c.z):
        if js(mix_dist(real, admitted + [(u,i)]),
              real_dist) <= eps:
            admitted.append((u, i))
\end{verbatim}

\section{补充实验：门控准入统计}
表~\ref{tab:admission} 报告完整框架各轮的候选数、准入数与实测
JS 距离。

\begin{table}[h]
\centering
\caption{门控准入统计（gated，种子 0）。}
\label{tab:admission}
\begin{tabular}{lcccc}
\toprule
轮次 & 1 & 2 & 3 & 4 \\
\midrule
候选数   & \AdmCandA & \AdmCandB & \AdmCandC & \AdmCandD \\
准入数   & \AdmPassA & \AdmPassB & \AdmPassC & \AdmPassD \\
实测 JS  & \AdmJsA & \AdmJsB & \AdmJsC & \AdmJsD \\
\bottomrule
\end{tabular}
\end{table}

\chapter*{致谢}
\addcontentsline{toc}{chapter}{致谢}
（按学校要求撰写：感谢导师的指导、实验室同学的帮助、家人的支
持等。）


\end{document}
